% ============================ Enrico Ribiani 16-03-2021 ====================================================================
% Base per i documenti  
\documentclass[12pt]{article}
% ------------ pacchetti necessari ----------------
\usepackage[a4paper, total={6in, 8in},margin=1in]{geometry} % formattazione decente della pagina
\usepackage{graphicx}                            % need for figure
\usepackage{amsmath}
\usepackage{amsfonts}                            % if you want the fonts
\usepackage{amssymb}                             % if you want extra symbols
\usepackage{graphicx}  
\renewcommand{\figurename}{Figura}  
\renewcommand{\contentsname}{Indice}                        % need for figures
\usepackage{mathptmx}
\usepackage{float}                               % serve per mettere tabelle e immagini dove si vuole 
\usepackage[utf8]{inputenc}
\usepackage{textcomp}
\usepackage[hang,flushmargin,bottom]{footmisc}   % footnote format
\usepackage{fancyhdr, lastpage}
\usepackage{titlesec}
\usepackage[table,dvipsnames]{xcolor}
%\pagestyle{fancy}
%\renewcommand{\headrulewidth}{0pt}
%\renewcommand*\contentsname{Indice}
\titleformat{\section}{\normalsize\bfseries}{\thesection.}{1em}{}	% required for heading numbering style
\titleformat*{\section}{\Large\bfseries}
\titleformat*{\subsection}{\large\bfseries}
%\usepackage{siunitx}
%\usepackage{tikz}
\usepackage{circuitikz}
\usepackage{multicol}
%\usepackage[siunitx]{circuitikz}
\usepackage{multirow}
\usepackage{tikz}
\usepackage{amsmath}
\usetikzlibrary{angles,quotes}
\usepackage{placeins}

\usepackage{wasysym}
%===================links=================
\usepackage{hyperref}
\hypersetup{
    colorlinks=true,
    linkcolor=darkgray,
    filecolor=Green,      
    urlcolor=Cyan,
    pdftitle={SAMPLE},
    pdfpagemode=FullScreen,
    }
%===================inizio pagina del titolo=================
\begin{document}
    \begin{titlepage}
    \begin{center}
% ------------------ inizio immagine logo ----------
\begin{figure}
    \centering
    \includegraphics{~/varie/logo.png}
    \label{fig:logo}
\end{figure}
% ------------------ fine immagine logo ----------
% ------------------ fine immagine logo ----------
-------------------------------------------------------------------------------------\\
\vspace{2\baselineskip}
\large Prova n°1
\hfill
\large $5^a$   AUB\\
\begin{flushleft}
    \large Enrico Ribiani\\
    \large Daniel Graziadei\\
    \large Gruppo 11\\
\end{flushleft}


\vfill

\Huge{\textbf{Amplificatore invertente e non invertente}}\\
\vfill
\vfill
\large{dd-mm-yyyy}
\end{center}
%=============== fine pagina titolo ===============
\end{titlepage}
\tableofcontents
\newpage
\vskip 1cm
\section{Scopo}
%
\section{Funzionamento}

\section{Schema}

\section{Materiale e Strumenti}
\label{Materiale e Strumenti}
\begin{multicols}{2}
    \begin{itemize}
    \item test
    \item test 2
    \end{itemize}
    \vfill\null
    \columnbreak
    \begin{itemize}
    \item test 3
    \item test 4
    \item test 5
    \end{itemize}
    \vfill\null
    \end{multicols}
\vspace{15pt}
\section{Contenuti Teorici}

\subsection{Descrizione della Prova}

\section{Raccolta dei dati}

\subsection{Tabella}

\subsection{commento dei dati raccolti in tabella}

\section{Elaborazione dei dati raccolti}

\section{Analisi critica dei risultati e conclusioni}
\end{document}